\usepackage[left=1.50cm, right=2.00cm, top=2.00cm, bottom=2.00cm]{geometry} % Tùy chỉnh lề
\usepackage{mathpazo} % Gói dùng để thay đổi phông chữ mặc định của tài liệu sang phông Palatino, và cung cấp phông chữ tương thích cho các ký hiệu toán học
\usepackage{graphicx} % Gói để chèn và xử lý hình ảnh: \includegraphics[cấu hình]{tên file} 
\usepackage{xcolor} % Gói tùy chỉnh màu sắc: \textcolor{màu}{Văn bản}, \colorbox{màu}{Văn bản}, \fcolorbox{màu 1}{màu 2}{Văn bản}, \definecolor{tên màu mới}{mô hình màu sắc}{thông số màu}
\usepackage{fontspec}
\usepackage{tabularx}
\usepackage{caption}
\usepackage{tcolorbox}
\setmainfont{Times New Roman} % hoặc font có hỗ trợ tiếng Việt
%\usepackage{xeCJK} % nếu cần tiếng Việt/tiếng Á
% Font hỗ trợ Devanagari (tiếng Hindi)
\newfontfamily\hindi{Noto Sans Devanagari}[Script=Devanagari, Scale=MatchLowercase]
\usepackage{float}
\usepackage{listingsutf8} % Gói dùng để hiển thị mã nguồn (code) chuyên nghiệp và có tô màu cú pháp trong tài liệu
\usepackage{booktabs}
\setlength{\parskip}{0.8em}
%------Định nghĩa listings C++------ 
\definecolor{dkgreen}{rgb}{0,0.6,0}
\definecolor{mauve}{rgb}{0.58,0,0.82}
\lstset{ 
frame=single, language=C++, aboveskip=3mm, belowskip=3mm, showstringspaces=false, columns=flexible, basicstyle={\normalsize\ttfamily}, numbers=none, % Không hiển thị số dòng 
keywordstyle=\color{blue}, commentstyle=\color{dkgreen}, stringstyle=\color{mauve}, breaklines=true, breakatwhitespace=true, tabsize=3,
literate={đ}{{\dj}}1 {Đ}{{\DJ}}1
           {á}{{\'a}}1 {à}{{\`a}}1 {ả}{{\h a}}1 {ã}{{\~a}}1 {ạ}{{\d a}}1
           {ă}{{\u a}}1 {ắ}{{\'\u a}}1 {ằ}{{\`\u a}}1 {ẳ}{{\h\u a}}1 {ẵ}{{\~\u a}}1 {ặ}{{\d\u a}}1
           {â}{{\^a}}1 {ấ}{{\'\^a}}1 {ầ}{{\`\^a}}1 {ẩ}{{\h\^a}}1 {ẫ}{{\~\^a}}1 {ậ}{{\d\^a}}1
           {é}{{\'e}}1 {è}{{\`e}}1 {ẻ}{{\h e}}1 {ẽ}{{\~e}}1 {ẹ}{{\d e}}1
           {ê}{{\^e}}1 {ế}{{\'\^e}}1 {ề}{{\`\^e}}1 {ể}{{\h\^e}}1 {ễ}{{\~\^e}}1 {ệ}{{\d\^e}}1
           {í}{{\'i}}1 {ì}{{\`i}}1 {ỉ}{{\h i}}1 {ĩ}{{\~i}}1 {ị}{{\d i}}1
           {ó}{{\'o}}1 {ò}{{\`o}}1 {ỏ}{{\h o}}1 {õ}{{\~o}}1 {ọ}{{\d o}}1
           {ô}{{\^o}}1 {ố}{{\'\^o}}1 {ồ}{{\`\^o}}1 {ổ}{{\h\^o}}1 {ỗ}{{\~\^o}}1 {ộ}{{\d\^o}}1
           {ơ}{{\ohorn}}1 {ớ}{{\'\ohorn}}1 {ờ}{{\`\ohorn}}1 {ở}{{\h\ohorn}}1 {ỡ}{{\~\ohorn}}1 {ợ}{{\d\ohorn}}1
           {ú}{{\'u}}1 {ù}{{\`u}}1 {ủ}{{\h u}}1 {ũ}{{\~u}}1 {ụ}{{\d u}}1
           {ư}{{\uhorn}}1 {ứ}{{\'\uhorn}}1 {ừ}{{\`\uhorn}}1 {ử}{{\h\uhorn}}1 {ữ}{{\~\uhorn}}1 {ự}{{\d\uhorn}}1
           {ý}{{\'y}}1 {ỳ}{{\`y}}1 {ỷ}{{\h y}}1 {ỹ}{{\~y}}1 {ỵ}{{\d y}}1 ,
    backgroundcolor=\color{gray!10}, % �� nền trắng xám nhạt (gray!10 = 10% xám)
    rulecolor=\color{gray!90},       % khung viền màu xám nhẹ
} 
\renewcommand{\figurename}{Hình}
\renewcommand{\tablename}{Bảng}
\lstdefinestyle{numbered}{ % Style cho code có đánh số 
numbers=left, % Hiển thị số dòng bên trái 
numberstyle=\small\color{gray}, numbersep=5pt, % Khoảng cách giữa số dòng và mã nguồn 
frame=single, framexleftmargin=15pt, % Đẩy khung bao quanh để chứa số dòng 
xleftmargin=15pt, % Lề trái để số dòng không bị cắt 
language=C++, aboveskip=3mm, belowskip=3mm, showstringspaces=false, columns=flexible, basicstyle={\normalsize\ttfamily}, keywordstyle=\color{blue}, commentstyle=\color{dkgreen}, stringstyle=\color{mauve}, breaklines=true, breakatwhitespace=true, tabsize=3,
literate={đ}{{\dj}}1 {Đ}{{\DJ}}1
           {á}{{\'a}}1 {à}{{\`a}}1 {ả}{{\h a}}1 {ã}{{\~a}}1 {ạ}{{\d a}}1
           {ă}{{\u a}}1 {ắ}{{\'\u a}}1 {ằ}{{\`\u a}}1 {ẳ}{{\h\u a}}1 {ẵ}{{\~\u a}}1 {ặ}{{\d\u a}}1
           {â}{{\^a}}1 {ấ}{{\'\^a}}1 {ầ}{{\`\^a}}1 {ẩ}{{\h\^a}}1 {ẫ}{{\~\^a}}1 {ậ}{{\d\^a}}1
           {é}{{\'e}}1 {è}{{\`e}}1 {ẻ}{{\h e}}1 {ẽ}{{\~e}}1 {ẹ}{{\d e}}1
           {ê}{{\^e}}1 {ế}{{\'\^e}}1 {ề}{{\`\^e}}1 {ể}{{\h\^e}}1 {ễ}{{\~\^e}}1 {ệ}{{\d\^e}}1
           {í}{{\'i}}1 {ì}{{\`i}}1 {ỉ}{{\h i}}1 {ĩ}{{\~i}}1 {ị}{{\d i}}1
           {ó}{{\'o}}1 {ò}{{\`o}}1 {ỏ}{{\h o}}1 {õ}{{\~o}}1 {ọ}{{\d o}}1
           {ô}{{\^o}}1 {ố}{{\'\^o}}1 {ồ}{{\`\^o}}1 {ổ}{{\h\^o}}1 {ỗ}{{\~\^o}}1 {ộ}{{\d\^o}}1
           {ơ}{{\ohorn}}1 {ớ}{{\'\ohorn}}1 {ờ}{{\`\ohorn}}1 {ở}{{\h\ohorn}}1 {ỡ}{{\~\ohorn}}1 {ợ}{{\d\ohorn}}1
           {ú}{{\'u}}1 {ù}{{\`u}}1 {ủ}{{\h u}}1 {ũ}{{\~u}}1 {ụ}{{\d u}}1
           {ư}{{\uhorn}}1 {ứ}{{\'\uhorn}}1 {ừ}{{\`\uhorn}}1 {ử}{{\h\uhorn}}1 {ữ}{{\~\uhorn}}1 {ự}{{\d\uhorn}}1
           {ý}{{\'y}}1 {ỳ}{{\`y}}1 {ỷ}{{\h y}}1 {ỹ}{{\~y}}1 {ỵ}{{\d y}}1
}

\usepackage{fancyvrb} % Needed for verbatim and \FancyVerbGetLine
\usepackage{amsmath, amssymb} % Gói lệnh toán học
\usepackage{amsfonts} % Gói dùng để mở rộng các ký hiệu toán học, đặc biệt là các ký hiệu tập hợp và phông chữ toán học đặc biệt
\usepackage{array} % Gói lệnh thêm các chức năng mở rộng cho bảng 
\usepackage{hhline} % Gói lệnh cải tiến hline: - một gạch ngang, = hai gạch ngang, ~ trống gạch ngang, # hai gạch xuống cắt qua hai gạch ngang, | một gạch xuống
\usepackage{fancyhdr} % Gói dùng để tùy chỉnh header (đầu trang) và footer (chân trang) của tài liệu
\usepackage[hidelinks]{hyperref} % Gói tạo liên kết trong và ngoài tài liệu
\usepackage{tikz} % Gói để vẽ hình ảnh vector, sơ đồ, biểu đồ, hình học, hình minh họa, sơ đồ khối, cây phân cấp, đồ thị mạng,...
\usetikzlibrary{shapes.geometric, arrows} % Dùng để nạp thêm thư viện con cho TikZ, giúp vẽ các sơ đồ khối (flowchart) hoặc các hình học đặc biệt.
\definecolor{fitblue}{RGB}{0, 80, 150} % Mã màu gần giống với logo
\usepackage{tocloft}
\renewcommand{\cfttoctitlefont}{\LARGE\bfseries} % căn giữa và đổi kích thước
\setlength{\cftbeforetoctitleskip}{0pt}
\setlength{\cftaftertoctitleskip}{0pt}\usepackage{xcolor}     % Cho màu sắc
\usepackage{etoolbox}   % Cho phép chỉnh sửa môi trường LaTeX

\makeatletter
\renewcommand{\@cite}[2]{\textcolor{blue}{[#1\if@tempswa , #2\fi]}}
\makeatother

%------Gói lệnh hỗ trợ bảng------ 
\usepackage{multirow} % Gộp hàng
\usepackage{multicol} % Gộp cột
\usepackage{booktabs} % Tùy chỉnh bảng đẹp hơn
\usepackage{colortbl} % Thêm màu
\usepackage{longtable} % Kéo dài bảng qua các trang
\usepackage{fancybox}
\usepackage{setspace}
\usepackage{booktabs}
\usepackage{float}
%------Định nghĩa kiểu cột có khả năng tùy chỉnh kích thước-----
\newcolumntype{L}[1]{>{\raggedright\let\newline\\\arraybackslash\hspace{0pt}}m{#1}}
\newcolumntype{C}[1]{>{\centering\let\newline\\\arraybackslash\hspace{0pt}}m{#1}}
\newcolumntype{R}[1]{>{\raggedleft\let\newline\\\arraybackslash\hspace{0pt}}m{#1}}

% Cấu hình BibLaTeX (Sử dụng style IEEE cho đúng chuẩn kỹ thuật)
\usepackage[backend=biber, style=ieee]{biblatex}
\addbibresource{references.bib} % Trỏ đến file dữ liệu

% Tùy chỉnh màu sắc cho link trích dẫn (nếu muốn)
\hypersetup{colorlinks=true, citecolor=blue, urlcolor=fitblue}